\documentclass[10pt,a4paper]{amsart}
\usepackage[margin=19mm]{geometry}
\usepackage{amsmath,amssymb,mathtools}
\usepackage[scr=rsfs,cal=euler]{mathalfa}
\usepackage{xfrac}
\usepackage[unicode,hidelinks,bookmarks=false]{hyperref}
\newcommand{\ie}{i.e.\ }
\newcommand{\cf}{cf.\ }
\newcommand{\resp}{respectively\ }
\newcommand{\Subsection}{\subsection}
\providecommand{\nobreakdash}{\nobreak\hbox{-}}
\setlength{\emergencystretch}{2em}
\multlinegap0pt
\usepackage{bm}
\usepackage{bbm}
\usepackage{dsfont}
\usepackage{xspace}
\usepackage{cancel}
\usepackage{mathtools}
\usepackage{amsthm}
\makeatletter
\@ifundefined{thm}{\newtheorem{thm}{Thm}}{}
\@ifundefined{prop}{\newtheorem{prop}[thm]{Proposition}}{}
\makeatother
\title[Beyond water limitation in vegetation-autotoxicity patternin]{Beyond water limitation in vegetation-autotoxicity patterning: a cross-diffusion model}
\author{Francesco Giannino, Annalisa Iuorio, Cinzia Soresina}
\date{Source: arXiv:2506.03981v1 (CC BY 4.0)}
\begin{document}
\maketitle

\noindent\emph{Source and licence.} Beyond water limitation in vegetation-autotoxicity patterning: a cross-diffusion model. Version arXiv:2506.03981v1 (CC BY 4.0), distributed under the Creative Commons Attribution 4.0 International licence, as recorded in the arXiv licence metadata for this exact version and shown on the abstract page https://arxiv.org/abs/2506.03981v1. Full rights notice: licenses/arxiv-2506.03981-CC-BY-4.0.txt.\par\smallskip

\noindent\textit{Editorial scope.} Counted unit: the Prop whose source label is \texttt{prop:stabcross} in the article (source0.tex lines 745--752, source label \texttt{prop:stabcross}), delivered together with the article's own complete original proof of it (source0.tex lines 753--791, one \texttt{proof} environment of 39 lines). No mathematics is written, repaired or re-derived: statement and proof are the article's own text line for line. Required context retained uncounted: eq:mainsys (lines 494--501); eq:mainsyshom (lines 574--581); prop:stabhom (lines 700--702). Every definition, display and label of those lines is retained unchanged. Omitted from this package and not redistributed: 1--493, 502--573, 582--699, 703--744, 792--1198. Nothing inside a retained excerpt is omitted or altered: every retained body is delivered exactly as the article's lines are written, so no omission receipt of any kind accompanies this package. Citations: the retained excerpts carry no citation command at all, so no bibliography entry of the article is transcribed and no reference is redistributed. Reference adaptation: none was needed and none was made; every reference inside the delivered unit resolves inside it, so no unresolved reference or citation is produced. Printed slips kept literal: no printed word, symbol, hypothesis or number of the counted statement or of its proof was altered. Layout and class: the article's own class is \texttt{elsarticle} with packages including fontenc, amsfonts, amsthm, amsmath, amssymb, bm, wasysym, accents, fixme, epsfig, overpic, psfrag, titlesec, verbatim; the wrapper is the lane's pinned \texttt{amsart} editorial wrapper at 10pt on A4, which restates the theorem-like environments the retained text needs and adds the article's own macro definitions that the retained text needs (none), with extra wrapper packages (bm, bbm, dsfont, xspace, cancel, mathtools). The delivered numbering is the unit's own. Assets: no retained span contains a figure environment, a graphics-inclusion command, a PDF-page inclusion, a table or a TikZ picture; no image file is copied into this package, the package declares no asset dependency and no image file is redistributed with it. Licence: version arXiv:2506.03981v1 of this article is distributed under the Creative Commons Attribution 4.0 International licence, as recorded in the arXiv licence metadata for this exact version and shown on the abstract page https://arxiv.org/abs/2506.03981v1. A full rights notice is delivered at licenses/arxiv-2506.03981-CC-BY-4.0.txt. Characters: every retained excerpt is pure ASCII, so the delivered engine, XeLaTeX with its default Latin Modern text font, reports no missing character.
\begin{equation} \label{eq:mainsys}
\begin{cases}
\partial_t R&= \Delta_x \left(\left(d_R -\sigma \theta(T)\right)R\right)
+\left(g-\gamma \theta(T)\right)R\left(1-\dfrac{R}{\hat{R}}\right)
-\left(d+s\theta(T)\right)R,\\[0.3cm]
\partial_t T&= d_T \Delta_x T +c\left(d+s\theta(T)\right)R-kT,
\end{cases}
\end{equation}

\begin{equation} \label{eq:mainsyshom}
\begin{cases}
 \dot{R}&= 
\left(g-\gamma \theta(T)\right)R\left(1-\dfrac{R}{\hat{R}}\right)
-\left(d+s\theta(T)\right)R,\\[0.3cm]
\dot{T}&= c\left(d+s\theta(T)\right)R-kT,   
\end{cases}
\end{equation}

\begin{prop} \label{prop:stabhom}
    The equilibrium $E_0$ of Eq.~\eqref{eq:mainsyshom} is unstable, whereas $E_*$ is asymptotically stable for the parameter values supporting its existence.
\end{prop}

\begin{prop} \label{prop:stabcross}
    The necessary condition for the instability of the coexistence steady-state ${E_*=(R_*,T_*)}$ with respect to heterogeneous perturbations is
    \begin{equation} \label{eq:condTur}
    \sigma > \sigma_L:=
    \left(\dfrac{k\hat{T}-csR_*}{cdR_*+kT_*}\right)d_R 
    + \dfrac{R_*}{\hat{R}}\left(\dfrac{d\hat{T}+sT_*}{(1-R_*/\hat{R})(cdR_*+kT_*)} \right)d_T.
\end{equation}
\end{prop}

\begin{proof}
The characteristic matrix $M_\kappa$ related to the $\kappa$-th eigenvalue of the Laplace operator with homogeneous Neumann boundary conditions, associated to system \eqref{eq:mainsys}, and evaluated at $E_*$ can be written as
\begin{equation}
    M_\kappa = J(R_*,T_*) -\lambda_\kappa J_\Delta(R_*,T_*) = \begin{pmatrix} 
J_{11}^*-\lambda_\kappa J^{\Delta*}_{11} & J_{12}^*-\lambda_\kappa J^{\Delta*}_{12} \\[0.3cm] 
J_{21}^* & J_{22}^*-\lambda_\kappa d_T
\end{pmatrix},
\end{equation}
where $J(R_*,T_*)$ is the Jacobian matrix for the homogeneous system and $J_\Delta(R_*,T_*)$ is the linearised diffusion matrix (both evaluated at $E_*$). In particular, the latter is defined as
\begin{equation*}
J_\Delta(R_*,T_*)=
\begin{pmatrix} 
J^{\Delta*}_{11} & J^{\Delta*}_{12} \\[0.2cm]
J^{\Delta*}_{21} & J^{\Delta*}_{22} 
\end{pmatrix}=
\begin{pmatrix} 
d_R-\sigma \theta(T_*) & -\sigma R_* \theta'(T_*)\\[0.3cm] 
0 & d_T
\end{pmatrix}.
\end{equation*}
For this matrix we have $J^{\Delta*}_{11},\,J^{\Delta*}_{22}>0$ and $J^{\Delta*}_{12}<0$.\\
In view of the assumption $\sigma<d_R$ and the consideration that $T_*<\hat{T}$, the trace of $M_\kappa$ remains negative, since
\[
\textnormal{tr}M_\kappa=\textnormal{tr}J(R_*,T_*)-\lambda_\kappa\left(d_T+d_R-\sigma \theta(T_*)\right)<0.
\]
On the other hand, the determinant of $M_\kappa$ reads
\[
\det M_\kappa=
d_T J^{\Delta*}_{11} \lambda_\kappa^2
-\left( d_T J^{*}_{11}+ J^{\Delta*}_{11}J^{*}_{22}
-J^{\Delta *}_{12}J^{*}_{21}\right) \lambda_\kappa
+\det J(R_*,T_*).
\]
Since $d_T>0$ and $\det (R_*,T_*)>0$ (due to Prop.~\ref{prop:stabhom}), in order for $E_*$ to become unstable -- thus leading to Turing instability -- the following necessary condition must hold
\[
d_T J^{*}_{11}+ J^{\Delta*}_{11}J^{*}_{22} -J^{\Delta *}_{12}J^{*}_{21}>0.
\]
Expressing all the Jacobian components explicitly in terms of $R_*$, $T_*$, we obtain Eq.~\eqref{eq:condTur}.
\end{proof}

\end{document}
